%! Least-Squares Regression — Unit 2, Lesson 2.4 — Lesson Plan
% GRADUAL-RELEASE plan (warm-up / guided notes and practice / debrief / close and start the
% homework). Teacher-facing — use the formal AP vocabulary freely.
% THERE IS NO GROUP ACTIVITY in this course, and NO INDEPENDENT PRACTICE SET in the notes:
% the notes are one Main Ideas / Notes table: I do / we do row by row, then the Guided
% Practice row together, and stop. The HOMEWORK is the individual practice, started in class
% in the last 8 minutes. Every box is `skillbox` (breakable) with a \boxguard ahead of it; this
% course has no `fixedskillbox`, no tier boxes, and no exit ticket.
\documentclass[10pt]{article}
\usepackage{apstats-article}
\usepackage{apstats-boxes}

\newcommand{\UnitNumberName}{Unit 2: Exploring Two-Variable Data \quad}
\newcommand{\LessonNumberName}{Lesson 2.4: Least-Squares Regression}

\begin{document}
\begin{center}
    {\Large\bfseries \CourseName \\
    {\small\bfseries \UnitNumberName \LessonNumberName}}
\end{center}

\begin{tcolorbox}[colback=sky, colframe=navy, boxrule=0.9pt, arc=2mm,
    left=2mm, right=2mm, top=1.5mm, bottom=1.5mm]
    \textbf{Primary Objective:} Students will compute the least-squares regression line from
    summary statistics, read its slope, intercept, and $r^2$ from computer output, and interpret
    each in context --- the slope as a change in the \emph{predicted} response, the intercept
    only when $x = 0$ is sensible, and $r^2$ as the share of the variation in $y$ explained by
    the linear model.
    \\[2pt]
    \textbf{CED alignment:} Topic 2.8 \quad$\cdot$\quad Big Idea: DAT (Data Analysis)
    \quad$\cdot$\quad Skills 2.C, 4.B \quad$\cdot$\quad LO DAT-1.G, DAT-1.H \quad$\cdot$\quad
    EK DAT-1.G.1--1.G.4, DAT-1.H.1--1.H.3
    \\[2pt]
    \textbf{Lesson model:} \emph{Gradual release --- I do, we do, you do.} There is no group
    activity. The guided notes deliver the vocabulary and the core idea directly, row by row
    in one \emph{Main Ideas / Notes} table --- each idea a definition to complete and a grid of
    short problems, the first worked by you and the rest by the class; the last row, Guided
    Practice, is worked together with students holding the pen;
    \textbf{the homework is the individual practice}, and students start it in class while you
    circulate. The whole-class debrief is \emph{spoken}.
    \\[2pt]
    \textbf{Compact budget:} the packet is six pages --- cover, warm-up, two pages of notes (14
    problems), one page of AP Practice, one page of homework --- so the whole lesson finishes in
    one period. Every display (two regression outputs) is \textbf{given}; students compute and
    interpret, they never draw.
\end{tcolorbox}

\renewcommand{\arraystretch}{1.4}
\boxguard
\begin{skillbox}[Learning Targets \& Key Understandings]{goldbox}
\begin{tabularx}{\linewidth}{@{}X|X@{}}
\textbf{Learning targets (student language)} & \textbf{Key understandings (the ``why'')} \\[2pt]
\hline
\vspace{-6pt}
\begin{itemize}[leftmargin=1.1em, nosep, after=\vspace{2pt}]
  \item I can explain what makes the least-squares line the best line, and compute its slope
        and intercept from $\bar x$, $s_x$, $\bar y$, $s_y$, and $r$.
  \item I can find the slope, the intercept, and $r^2$ in computer regression output.
  \item I can interpret the slope and the $y$-intercept in context, and say when the intercept
        means nothing.
  \item I can interpret $r^2$ in context.
\end{itemize}
&
\vspace{-6pt}
\begin{itemize}[leftmargin=1.1em, nosep, after=\vspace{2pt}]
  \item The LSRL minimizes the sum of the squared residuals and passes through
        $(\bar x, \bar y)$; $b = r\,s_y/s_x$ and $a = \bar y - b\bar x$ (EK DAT-1.G.1, 1.G.2,
        1.H.3).
  \item \textbf{Target misconception:} reading the slope as a promise about every individual
        --- ``when a student sleeps one more hour, their reaction time goes down 24~ms.'' The
        slope is the change in the \emph{predicted} $y$ per unit of $x$ (EK DAT-1.H.2); real
        points scatter around the line, and observational data show no cause.
  \item \textbf{Second misconception:} interpreting the intercept when $x = 0$ is outside the
        data or impossible (EK DAT-1.G.3). Sleep: meaningless. Used Civics: age 0 is a new car,
        so it means something.
  \item \textbf{Third misconception:} taking $r = +\sqrt{r^2}$. The sign of $r$ is the sign of
        the slope; $r^2$ is a proportion of \emph{variation} explained (EK DAT-1.G.4).
\end{itemize}
\end{tabularx}
\end{skillbox}

\boxguard
\begin{skillbox}[{Vocabulary, Concepts, \& Theorems --- taught directly in the guided notes}]{greenbox}
\begin{minipage}[t]{0.48\linewidth}
    \begin{itemize}
        \item \textbf{Least-squares regression line (LSRL)} --- minimizes $\sum(y - \hat y)^2$;
              contains $(\bar x, \bar y)$ (EK DAT-1.G.1)
        \item \textbf{Slope} $b = r\,\dfrac{s_y}{s_x}$ --- the change in predicted $y$ for each
              1-unit increase in $x$ (EK DAT-1.G.2, 1.H.2)
        \item \textbf{Intercept} $a = \bar y - b\bar x$ --- predicted $y$ at $x = 0$; sometimes
              no logical meaning (EK DAT-1.G.3, 1.H.3)
    \end{itemize}
\end{minipage}
\hfill
\begin{minipage}[t]{0.48\linewidth}
    \begin{itemize}
        \item \textbf{Coefficient of determination} $r^2$ --- the proportion of the variation in
              $y$ explained by the linear model with $x$ (EK DAT-1.G.4)
        \item \textbf{Computer output} --- Constant row $= a$; the $x$ row $= b$; R-Sq $= r^2$
              (EK DAT-1.H.1). $S$ is printed but not taught today.
        \item \textbf{Carried forward} --- $r$ (2.2); $\hat y = a + bx$, residual
              $= y - \hat y$, extrapolation (2.3)
    \end{itemize}
\end{minipage}
\end{skillbox}

% A tabularx never splits, so guard generously ahead of it.
\boxguard[30]
\begin{skillbox}[Lesson at a Glance --- \MeetingLength]{sky}
\renewcommand{\arraystretch}{1.25}
\begin{tabularx}{\linewidth}{@{}l c X X@{}}
\textbf{Phase} & \textbf{Min} & \textbf{Students} & \textbf{Teacher} \\
\hline
Warm-Up & 5 & A residual, two lines' residuals summed and squared, and the arithmetic of
$b = r\,s_y/s_x$ & Debrief item~2 aloud; leave ``both sums are 0'' on the board \\
Guided Notes \& Practice & 34 & Fill the vocabulary box and problems 1--10 row by row, then
\emph{Used Cars} (11--14) together, pen in their hand & \textbf{I do} each row's definition lines
and problems 1, 4, 7; \textbf{we do} the rest (24) $\to$ \textbf{we do} Guided Practice (10) \\
Debrief & 8 & Pens down --- answer aloud, nothing to fill in & Open on an \emph{almost}-right
problem~10; then the two intercepts side by side \\
Close \& Start the Homework & 8 & \textbf{Start the homework in class}, alone and silent &
Circulate for the second formative read (A2); preview Lesson~2.5 \\
\end{tabularx}
\end{skillbox}

\boxguard[20]
\begin{skillbox}[Warm-Up --- Activate Prior Knowledge (5 min)]{sky}
\begin{minipage}[t]{0.48\linewidth}
    \textbf{The three items and what each seeds}
    \begin{itemize}
        \item \textbf{1.} Predict, then a residual. \textbf{Spirals} to Lesson 2.3: residual
              $=$ actual $-$ predicted $= 16 - 14 = 2$.
        \item \textbf{2.} Two lines, both residual sums 0, squares 20 and 16. \textbf{Seeds} the
              least-squares criterion in the first row of the notes.
        \item \textbf{3.} $(-0.80)^2 = 0.64$ and $-0.80 \times 45/1.5 = -24$. \textbf{Seeds}
              problems 1 and 5 --- the numbers are the notes' numbers, before they have names.
    \end{itemize}
\end{minipage}
\hfill
\begin{minipage}[t]{0.48\linewidth}
    \textbf{Running it}
    \begin{itemize}
        \item \textbf{Debrief item~2 aloud.} Ask why the plain sum cannot pick the better
              line; write ``both sums are 0'' on the board and leave it up --- the first row
              of the notes points back at it.
        \item Do not name $r^2$ or the slope formula on item~3. Just say ``remember these two
              numbers'' --- they reappear in problems 1 and 5 within fifteen minutes.
        \item Item~1 is a 60-second spiral. If the room is quick, move on.
    \end{itemize}
\end{minipage}
\end{skillbox}

\boxguard
\begin{skillbox}[Guided Notes \& Practice --- I do, then we do (34 min)]{sky}
\textbf{This is the whole lesson.} The notes are one \emph{Main Ideas / Notes} table --- three
instruction rows and a Guided Practice row, 14 short problems. \textbf{In every instruction row
you fill the definition lines and work the first problem (1, 4, 7); the class works the rest,
pen in their hand.} The vocabulary box is filled \emph{as you go}, never in advance.
\begin{multicols}{2}
\textbf{Least Squares (7 min).} Point at the warm-up's ``both sums are 0'': that is why we
square. Fill the two blanks, then the three steps. Work problem~1 yourself ($b = -24$); let the
class do 2 and 3. Insist problem~3 uses the \emph{variable names} ---
$\widehat{\text{time}} = 482.8 - 24(\text{sleep})$ --- not $x$ and $y$.

\textbf{Computer Output \& R-Squared (7 min).} Same line, now as software prints it. Do
problem~4 (find $a$ and $b$) yourself and point out that they match 1 and 2. \textbf{Problem~5
is the trap}: R-Sq $= 64\%$, so $r = \pm 0.80$ --- take a hand count before anyone speaks; the
room that says $+0.80$ forgot the negative slope. Problem~6 sets the $r^2$ sentence they will
reuse all unit.

\textbf{Slope \& Intercept (10 min) --- the crux.} Fill the definitions and the template box,
then work problem~7 aloud, stressing \emph{predicted}. Problem~8 is the intercept that means
nothing (0 hours is outside 4--10). Problem~9 makes them use the slope as a difference of
predictions. \textbf{Problem~10 is the crux}: Maya's sentence has the right number and the wrong
claim. Ask \emph{``Did every student who slept an extra hour get exactly 24 ms faster?''}, then
\emph{``Who decided how much each student slept?''}

\columnbreak
\textbf{We do (about 10 min): Used Cars.} A fresh context, fresh output, problems 11--14,
worked entirely together. Students hold the pen; you hold the questions: \emph{``Which row is
the slope?''} \quad \emph{``Is a 0-year-old Civic a real car?''} \quad \emph{``Which sign does
$r$ get?''} \textbf{Problem~12 is the crux transferred} --- if a student writes ``the price
drops \$1,850,'' make the room supply the missing word. \textbf{Problem~13 is the contrast}:
this intercept \emph{does} mean something, so students who learned ``intercepts are always
meaningless'' get caught here.

Circulate and demand the reason, not the number:
\begin{itemize}[leftmargin=1.1em, nosep]
  \item \textbf{Problem~11:} ``Where did the 24.60 come from --- which row?''
  \item \textbf{Problem~12:} ``Every car? Or the prediction?''
  \item \textbf{Problem~14:} ``$\sqrt{0.81}$ is 0.90. Is $r$ positive?''
\end{itemize}
While circulating, pick the papers for the debrief --- one \emph{almost}-right problem~10 or 12
(right number, no ``predicted''), one fully correct. \textbf{Formative read, and reteach
trigger:} if more than a third of the room writes problem~12 without ``predicted,'' reteach the
template box at the start of Lesson~2.5. Your second read comes in the last eight minutes, on
homework A2.
\end{multicols}
\end{skillbox}

\boxguard
\begin{skillbox}[Debrief --- whole class, spoken (8 min)]{sky}
Pens down, papers up. \textbf{This debrief is spoken} --- there is nothing on the page to fill in.
\begin{multicols}{2}
\begin{enumerate}[leftmargin=1.3em, nosep, itemsep=3pt]
  \item \textbf{Open on the \emph{almost}-right problem~10 (or 12)} you picked: right number,
        missing ``predicted.'' Ask the room what is true and what is overclaimed. Land two
        things --- the line describes \emph{predictions}, and data we only observed cannot show
        cause.
  \item \textbf{Two intercepts side by side:} 482.8~ms at 0 hours of sleep (no meaning ---
        outside the data) and \$24{,}600 at age 0 (a new Civic --- meaningful). The rule is
        not ``intercepts are meaningless''; it is ``ask whether $x = 0$ makes sense.''
  \item \textbf{The trap:} $r$ from R-Sq. The square root forgets the sign; the slope remembers
        it.
\end{enumerate}
\columnbreak
\textbf{Demand one sentence aloud:} a student states $r^2$ for the Civics in the template ---
``About 81\% of the variation in price is explained by the linear model with age.'' Reject
``81\% of the points are on the line'' and ``the model is 81\% accurate'' out loud.

\textbf{If time is short}, cut item~3 --- item~1 is the one that cannot be cut. \textbf{Do not
let the debrief eat the last eight minutes}; the homework start is the second formative read.
\end{multicols}
\end{skillbox}

\boxguard
\begin{skillbox}[AP Practice --- assigned, not scored]{goldbox}
A fresh context (a caf\'e's daily high temperature and iced drinks sold, with computer output).
\textbf{MC~1} is the crux in AP clothing --- the correct slope interpretation (B) against a
deterministic (A), a causal (C), an $r^2$ (D), and an intercept (E) distractor. \textbf{MC~2}
is the trap: $r = +0.850$ (D) from R-Sq $= 72.3\%$ and a positive slope, with $-0.850$ waiting
for students who guess the sign. \textbf{FRQ} (a) interpret $r^2$; (b) the intercept,
$-42.5$ drinks at 0$^\circ$F --- impossible and far outside 55--95$^\circ$F; \textbf{(c) is the
spiral} to Lesson 2.3: $\hat y = 145.5$, residual $= +4.5$ drinks. \textbf{Not scored --- the
cover reads NA.} Go over it at the start of Lesson~2.5.
\end{skillbox}

\boxguard
\begin{skillbox}[Homework --- scored, due the first class after two study halls]{goldbox}
Five items on car weight and highway mileage (a third context): \textbf{A1} compute $b = -7.0$
and $a = 54.5$ and write the line; \textbf{A2} interpret the slope; \textbf{A3} the intercept
(a 0-pound car --- meaningless); \textbf{B1} $r^2 \approx 0.71$ in context; \textbf{B2} an
AP-style multiple choice (B) with a required justification. \textbf{Scoring:} 10 points, 2 per
item. \textbf{A2 predicts next-lesson performance}: a student who drops ``predicted'' here has
not left the misconception behind. \textbf{Paper is the default.} \textbf{DeltaMath override
(occasional):} when DeltaMath carries this topic, assign it instead, strike the score cell on
the cover, and leave this page as optional practice --- never assign both.
\end{skillbox}

\boxguard
\begin{skillbox}[Watch For (while circulating)]{redbox}
\begin{itemize}[leftmargin=1.2em, nosep]
  \item \textbf{The slope read as a promise about every individual} (notes~10, 12; AP~MC1;
        HW~A2, B2 --- the target misconception). ``Their reaction time goes down 24 ms.'' Probe:
        \emph{``Every student? Or the prediction?''}
  \item \textbf{Cause read into the slope} (notes~10; AP~MC1(C); HW~B2(E)). Probe: \emph{``Who
        assigned the sleep?''}
  \item \textbf{Intercept interpreted outside the data --- or dismissed when it is sensible}
        (notes~8, 13; AP~(b); HW~A3). Probe: \emph{``Is $x = 0$ a real case here? Is it near
        the data?''}
  \item \textbf{$r = +\sqrt{r^2}$} (notes~5, 14; AP~MC2). Probe: \emph{``Which way does the
        line go?''}
  \item \textbf{$r^2$ as ``percent of points on the line'' or ``percent accurate''} (notes~6;
        HW~B1). Probe: \emph{``Percent of \emph{what}?''}
  \item \textbf{$x$ and $y$ swapped in $b = r\,s_y/s_x$} (notes~1; HW~A1). Probe:
        \emph{``Which variable is the response? Its $s$ goes on top.''}
\end{itemize}
Cold-call prompts: \emph{``Say the slope sentence without the word `predicted' --- now tell me
what is wrong with it.''} \quad \emph{``Name an $x$ where the intercept would mean something.''}
\end{skillbox}

\boxguard
\begin{skillbox}[Close \& Start the Homework (8 min)]{goldbox}
\textbf{Students start the homework here, in class, alone and silent --- this is the individual
practice.} Hand the launch first --- five items on paper, scored, due the first class after two
study halls --- and point out that the AP Practice page before it is theirs to work and bring to
review. Put students on A1 and A2 while you can still catch a wrong start. \textbf{A2 is the
column that tells you whether the crux landed.} Sort as you walk: said ``predicted'' in context
(ready); right number, no ``predicted'' (ready, needs the language); wrote a promise or a cause
(the misconception survived) --- and open Lesson~2.5 by reteaching the template box if that
third pile runs past a third of the room. Then close on what changed: \emph{``You came in
thinking a line tells you what will happen. You leave knowing it tells you what to
\emph{predict} --- and how much of the variation it explains.''} \textbf{Preview:}
Lesson~2.5, \emph{Departures from Linearity} --- what one unusual point can do to $b$, $a$, and
$r^2$.
\end{skillbox}

% ── Teacher notes — one per component, in packet order (four: there is no activity) ──
\begin{teachernote}[Warm-Up]
Four minutes. Item~2 is the seed: both residual sums are 0, so the plain sum is useless and
squaring is what separates the lines --- say it once and leave ``both sums are 0'' on the board.
Item~3's numbers ($0.64$ and $-24$) are the notes' numbers on purpose; do not name them yet.
Item~1 is a 60-second spiral to Lesson~2.3 --- do not let it expand.
\end{teachernote}

\begin{teachernote}[Guided Notes \& Practice]
\textbf{Pace it 24 / 10, then 8 for the debrief, then 8 for the homework start.} Rows 1 and 2 are
machinery and must move --- seven minutes each. You work problems 1, 4, and 7; the class works
the rest of each grid. \textbf{Problem~5 is the trap} (the sign of $r$), and \textbf{problem~10
is the crux}: Maya's sentence has the right number and claims too much. Do not resolve it too
fast --- let a student defend it, then ask whether \emph{every} student who slept an extra hour
got exactly 24~ms faster. The fix is one word, \emph{predicted}, plus the reminder that nobody
assigned the sleep.

\textbf{Used Cars is where the pen changes hands.} Problem~12 is the crux transferred;
problem~13 is the contrast --- an intercept that \emph{does} mean something. If you only get
through 11, 12, and 13, that is the right trade. The output prints $S$; if a student asks, it is
the typical size of a residual --- it is not tested in this lesson.

\textbf{Debrief (8 min), spoken.} Open on the almost-right problem~10 or 12. \textbf{Do not let
the debrief eat the last eight minutes} --- the homework start is your second formative read.
\end{teachernote}

\begin{teachernote}[AP Practice]
\textbf{This page is not required and not scored} --- the graded practice is the homework behind
it. Go over it at the start of Lesson~2.5. \textbf{MC~2 is the one worth reading closely}:
students who choose (A) $-0.850$ are guessing the sign, and students who choose (C) $0.723$
have confused $r$ with $r^2$. On the free response, part (b) is the highest-value part: a
negative number of drinks is the cleanest possible case of an intercept with no meaning, and
students must say \emph{why} (impossible, and 0$^\circ$F is far outside the data).
\end{teachernote}

\begin{teachernote}[Homework]
\textbf{Scored, 10 points (2 per item), due the first class after two study halls.} The eight
minutes at the end of the period are a \emph{start}: A1 and A2 while you can still catch a wrong
start; A3, B1, and B2 go home. Read A2 first when scoring --- ``predicted'' in context earns the
2; a correct number without it earns 1. On A1, a student with $b = -0.84 \cdot 0.6/5$ swapped
the standard deviations; that is worth a one-line comment, not a reteach. B2 requires the
justification --- a bare (B) earns 1 point. If more than a third miss A2, reopen the template
box in the Lesson~2.5 warm-up debrief. \textbf{Paper is the default.} On the occasions you assign
DeltaMath in its place, strike the score cell on the cover and tell students this page is now
optional practice. Do not assign both.
\end{teachernote}
\end{document}
